\section{Social Conformity}
\label{sec:conformity}

\subsection{Definition}
\label{sec:conformity_definition}

Conformity is the adoption of a wrong answer to match a peer or majority position,
in an Asch-style paradigm adapted for language models: the model answers privately,
observes one or more peer answers, and is given the opportunity to revise. It
overlaps with epistemic capitulation (\S\ref{sec:epistemic_capitulation}) but is
defined by the presence of an attributed source -- one or more other agents, real
or simulated -- rather than bare disagreement.

\subsection{How It Is Recreated}
\label{sec:conformity_elicitation}

The dominant protocol constructs a simulated multi-agent setting: peer answers are
inserted with varying majority size, unanimity, framing (neutral, confident,
uncertain), and channel (labeled as another user, another assistant, or a tool
output), and the subject model's revision is scored. A methodologically important
control removes the attributed speaker entirely, replacing ``the other three
participants said X'' with a bare repeated assertion of X. This control is
essential: without it, a study cannot distinguish genuine social influence from a
purely textual repetition effect.

\subsection{Evidence and Analysis}
\label{sec:conformity_evidence}

Table~\ref{tab:conformity_evidence} summarizes representative results. The central
finding of this section is that a large majority of what the literature has called
``conformity'' does not require a social speaker at all: a bare repeated assertion
with no attributed source reproduces most of the effect obtained with a full
multi-agent framing, and adding an explicit source or authority label contributes
only a modest increment on top of that floor. This does not make conformity a null
effect -- the floor itself is large -- but it means that many existing conformity
benchmarks conflate a general susceptibility to repeated text with a genuinely
social phenomenon, and future work should report both quantities separately.
Conformity also correlates with the model's own uncertainty: models conform more
often when their initial confidence is low, consistent with an informational
influence account rather than a purely normative one, though the two are not
mutually exclusive and can co-occur in the same model.

\begin{table}[t]
\centering
\caption{Reported social conformity effects across studies.}
\label{tab:conformity_evidence}
\begin{tabular}{p{2.6cm}p{2.6cm}p{2.6cm}p{5.2cm}}
\toprule
\textbf{Study} & \textbf{Pressure Type} & \textbf{Model(s)} & \textbf{Effect Size} \\
\midrule
BenchForm & peer majority, Trust/Doubt priming & 9 models & Doubt-protocol accuracy drop up to 38.6\% \\
Not Just RLHF & multi-agent peer disagreement & 4 families (base + instruct) & yield 44--98\%; base $\geq$ instruct in 10/12 cells \\
Most Conformity Needs No Speaker & bare repeated assertion vs.\ attributed source & 6 open-weight models & 66.5\% floor with no speaker; +12.9pp for expert-panel framing \\
MUSE & distractor + authority framing & 7 models & conformity 29.0\% (certain) vs.\ 39.6\% (uncertain) \\
KAIROS & multi-agent support/opposition, rapport level & wide range & avg.\ $\leq$32B models $-$5.65\%, $>$32B $-$3.64\% \\
Human-like Social Compliance & single-authority + 3-agent consensus & 15 models & sycophancy--conformity correlation $R^2=0.31$ \\
\bottomrule
\end{tabular}
\end{table}

\subsection{Mitigation}
\label{sec:conformity_mitigation}

Introducing a single dissenting agent who argues correctly reduces conformity
substantially and generalizes across framings, more reliably than persona-based or
reflection-based prompting. Reinforcement learning against a multi-agent-context
reward has produced the only mitigation in this survey shown to improve both
robustness to incorrect peers and appropriate uptake of correct peer information
simultaneously, though the gain is concentrated in larger models. Chain-of-thought
prompting and generic self-reflection, by contrast, have been shown \emph{not} to
selectively reduce harmful conformity while preserving beneficial revision --
several studies find these generic reasoning interventions move both quantities
together rather than disentangling them.
